\chapter{Exact limits, reciprocity, and research frontiers}\label{chap:limits-reciprocity-frontiers}

This chapter isolates three principles that are easy to use informally and
equally easy to overstate: modular acceleration near \(q=1\), passage from
finite identities to formal or analytic limits, and stabilization at the two
ends of a reciprocal polynomial.  Each proved statement below records the
topology in which its limit is taken.  The final section then lists open
programs as problems rather than silently promoting them to conjectures.

\section{An exact modular acceleration}

The eta transformation in \Cref{qg:thm-theta-eta-modular} gives more than the
leading singular term of \((q;q)_\infty\).  It replaces a slowly convergent
product with one whose nome is exponentially small.

\begin{theorem}[Exact modular form of the Euler product]
\label{qg:thm-qpochhammer-modular-asymptotic}
Let \(t>0\), and put
\[
 q=\EulerE^{-t},
 \qquad
 Q=\EulerE^{-4\pi^2/t}.
\]
Then
\begin{equation}
 (q;q)_\infty
 =\sqrt{\frac{2\pi}{t}}
   \exp\!\left(-\frac{\pi^2}{6t}+\frac{t}{24}\right)
   (Q;Q)_\infty.
\label{eq:qg-qpochhammer-modular-exact}
\end{equation}
Consequently, as \(t\to0^+\),
\begin{align}
 (q;q)_\infty
 &=\sqrt{\frac{2\pi}{t}}
   \exp\!\left(-\frac{\pi^2}{6t}+\frac{t}{24}\right)
   \left(1+O\!\left(\EulerE^{-4\pi^2/t}\right)\right),
\label{eq:qg-qpochhammer-modular-asymptotic}\\
 \log(q;q)_\infty
 &=-\frac{\pi^2}{6t}
   +\frac12\log\frac{2\pi}{t}
   +\frac{t}{24}
   +O\!\left(\EulerE^{-4\pi^2/t}\right).
\label{eq:qg-qpochhammer-log-asymptotic}
\end{align}
The error terms are real and uniform for \(0<t\le t_0\), for every fixed
\(t_0>0\) small enough.
\end{theorem}

\begin{proof}
Set \(\tau=\ImaginaryUnit t/(2\pi)\).  The definition
\eqref{eq:qg-eta-definition} gives
\begin{equation}
 \eta\!\left(\frac{\ImaginaryUnit t}{2\pi}\right)
 =\EulerE^{-t/24}(\EulerE^{-t};\EulerE^{-t})_\infty.
\label{eq:qg-eta-at-small-imaginary-argument}
\end{equation}
Since \(-1/\tau=2\pi\ImaginaryUnit/t\), the inversion formula
\eqref{eq:qg-eta-S}, with its positive square root on the positive imaginary
axis, reads
\[
 \eta\!\left(\frac{\ImaginaryUnit t}{2\pi}\right)
 =\sqrt{\frac{2\pi}{t}}
   \eta\!\left(\frac{2\pi\ImaginaryUnit}{t}\right).
\]
Applying the definition once more at \(2\pi\ImaginaryUnit/t\) yields
\[
 \eta\!\left(\frac{2\pi\ImaginaryUnit}{t}\right)
 =\EulerE^{-\pi^2/(6t)}
   (\EulerE^{-4\pi^2/t};\EulerE^{-4\pi^2/t})_\infty.
\]
Combining the last two identities with
\eqref{eq:qg-eta-at-small-imaginary-argument} proves the exact formula
\eqref{eq:qg-qpochhammer-modular-exact}.

It remains to justify the stated exponential remainder rather than merely
observe that \(Q\to0\).  For \(0<Q\le1/2\), the elementary inequality
\(-\log(1-u)\le u/(1-u)\), \(0\le u<1\), gives
\[
 0\le-\log(Q;Q)_\infty
 =\sum_{m\ge1}-\log(1-Q^m)
 \le\sum_{m\ge1}\frac{Q^m}{1-Q^m}
 \le\frac{Q}{(1-Q)^2}
 \le4Q.
\]
Thus \(\log(Q;Q)_\infty=O(Q)\).  If
\(L=-\log(Q;Q)_\infty\), then \(0\le L\le4Q\) and
\[
 \bigl|(Q;Q)_\infty-1\bigr|=1-\EulerE^{-L}\le L\le4Q.
\]
Substitution in the exact formula gives the first asymptotic, and taking the
real logarithm gives the second.  The displayed inequalities also provide
the asserted uniformity.
\end{proof}

\section{Coefficientwise and \(q\)-adic limits}

For a series \(G(q)=\sum_{d\ge0}g_dq^d\), write
\(\operatorname{ord}_qG=\min\{d:g_d\ne0\}\), with
\(\operatorname{ord}_q0=\infty\).  The congruence
\(G\equiv H\pmod{q^{N+1}}\) means that their coefficients agree through
degree \(N\).  These conventions make the finiteness behind the usual
“truncate and pass to the limit” argument explicit.

\begin{proposition}[Coefficientwise passage to an analytic limit]
\label{qg:prop-coefficientwise-limit}
Let \(R\) be a commutative ring and let \(P_N\in R[q]\).  Suppose that for
every \(d\ge0\) there are \(c_d\in R\) and \(N_d\) such that
\[
 [q^d]P_N=c_d\qquad(N\ge N_d).
\]
Then there is a unique series \(F=\sum_{d\ge0}c_dq^d\in R[[q]]\), and
\(P_N\to F\) coefficientwise.

If \(R=\ComplexNumbers\) and the polynomials \(P_N\), viewed as holomorphic
functions, converge locally uniformly on the unit disk to a function \(P\),
then \(P\) is holomorphic and
\begin{equation}
 c_d=\frac{1}{2\pi\ImaginaryUnit}
     \int_{|z|=r}\frac{P(z)}{z^{d+1}}\,dz
     =\frac{P^{(d)}(0)}{d!}
 \qquad(0<r<1).
\label{eq:qg-coefficientwise-cauchy}
\end{equation}
In particular, the formal coefficientwise limit is precisely the Taylor
series of the analytic limit.
\end{proposition}

\begin{proof}
The first assertion is coefficientwise: the eventual value \(c_d\) determines
the coefficient of degree \(d\), so the displayed series exists and is the
only possible coefficientwise limit.

For the analytic assertion, local uniform convergence of holomorphic
functions makes \(P\) holomorphic.  Fix \(0<r<1\).  Cauchy's coefficient
formula and uniform convergence on the circle \(|z|=r\) give
\[
 [q^d]P_N
 =\frac{1}{2\pi\ImaginaryUnit}
   \int_{|z|=r}\frac{P_N(z)}{z^{d+1}}\,dz
 \longrightarrow
 \frac{1}{2\pi\ImaginaryUnit}
   \int_{|z|=r}\frac{P(z)}{z^{d+1}}\,dz.
\]
The left side is eventually equal to \(c_d\).  A second application of
Cauchy's formula identifies the last integral with \(P^{(d)}(0)/d!\), proving
\eqref{eq:qg-coefficientwise-cauchy}.
\end{proof}

\begin{proposition}[Exact truncation principle]
\label{qg:prop-truncation-principle}
Let \((e_j)_{j\ge1}\) be an arbitrary sequence of integers, and interpret
negative powers by their binomial expansions in \(\IntegerNumbers[[q]]\).
The finite products
\[
 F_M(q)=\prod_{j=1}^{M}(1-q^j)^{e_j}
\]
have a well-defined coefficientwise limit \(F\in\IntegerNumbers[[q]]\), and
for every \(N\ge0\),
\begin{equation}
 F(q)\equiv\prod_{j=1}^{N}(1-q^j)^{e_j}
 \pmod{q^{N+1}}.
\label{eq:qg-truncation-product}
\end{equation}

More generally, let \(R\) be a commutative ring and let
\((S_m)_{m\ge0}\subset R[[q]]\) satisfy
\(\operatorname{ord}_qS_m\to\infty\), in the precise sense that for each
\(N\) only finitely many \(m\) have
\(\operatorname{ord}_qS_m\le N\).  Then the \(q\)-adic sum
\(S=\sum_{m\ge0}S_m\) is well defined and
\begin{equation}
 S(q)\equiv
 \sum_{\operatorname{ord}_qS_m\le N}S_m(q)
 \pmod{q^{N+1}}.
\label{eq:qg-truncation-series}
\end{equation}
\end{proposition}

\begin{proof}
For every \(j\ge1\) and \(e\in\IntegerNumbers\),
\[
 (1-q^j)^e\equiv1\pmod{q^j}.
\]
For \(e\ge0\) this follows from the finite binomial theorem.  For
\(e=-r<0\), it follows from
\[
 (1-q^j)^{-r}
 =\sum_{k\ge0}\binom{k+r-1}{r-1}q^{jk}
 \qquad(r\ge1).
\]
Consequently, if \(M\ge N\), every factor with \(j>N\) is congruent to \(1\)
modulo \(q^{N+1}\), and hence
\[
 F_M(q)\equiv F_N(q)\pmod{q^{N+1}}.
\]
Thus every coefficient of \(F_M\) stabilizes; apply
\Cref{qg:prop-coefficientwise-limit} with \(R=\IntegerNumbers\).  Passing to
the stabilized coefficients gives \eqref{eq:qg-truncation-product}.

For the series assertion, fix a degree \(d\).  Only finitely many \(S_m\)
have order at most \(d\), so the coefficient
\(\sum_m[q^d]S_m\) is a finite sum in \(R\).  These finite sums define a
unique element \(S\in R[[q]]\).  Terms with order greater than \(N\) vanish
modulo \(q^{N+1}\), and the remaining index set is finite by hypothesis.
This proves \eqref{eq:qg-truncation-series}.
\end{proof}

\section{Borwein dissections and two-ended stabilization}

For \(n\ge1\), define
\begin{equation}
 P_n(q)=\frac{(q;q)_{3n}}{(q^3;q^3)_n}
 =\prod_{j=1}^{n}(1-q^{3j-2})(1-q^{3j-1}).
\label{eq:qg-borwein-product}
\end{equation}
Grouping coefficients by their exponent modulo \(3\) gives unique
polynomials \(A_n,B_n,C_n\in\IntegerNumbers[x]\) such that
\begin{equation}
 P_n(q)=A_n(q^3)-qB_n(q^3)-q^2C_n(q^3).
\label{eq:qg-borwein-dissection}
\end{equation}
The minus signs are part of the normalization; none of the arguments below
assumes a sign for the coefficients of the three polynomials.

\begin{proposition}[Reciprocity of the Borwein dissection]
\label{qg:prop-borwein-reciprocity}
For every \(n\ge1\),
\begin{align}
 A_n(x)&=x^{n^2}A_n(x^{-1}),\\
 B_n(x)&=x^{n^2-1}C_n(x^{-1}),\\
 C_n(x)&=x^{n^2-1}B_n(x^{-1}).
\label{eq:qg-borwein-reciprocity}
\end{align}
In particular,
\[
 B_n(1)=C_n(1),
 \qquad
 A_n(1)=2B_n(1).
\]
\end{proposition}

\begin{proof}
The degree of \(P_n\) is
\[
 \sum_{j=1}^{n}\bigl((3j-2)+(3j-1)\bigr)=3n^2.
\]
For positive integers \(a,b\),
\((1-q^{-a})(1-q^{-b})=q^{-a-b}(1-q^a)(1-q^b)\).
There are two factors in every \(j\)-block, so their signs cancel, and
\[
 P_n(q)=q^{3n^2}P_n(q^{-1}).
\]
Insert \eqref{eq:qg-borwein-dissection} on both sides and put \(x=q^3\).
The residue-zero terms give
\(A_n(x)=x^{n^2}A_n(x^{-1})\).  The residue-one terms on the left match the
reflected residue-two terms on the right, giving
\(B_n(x)=x^{n^2-1}C_n(x^{-1})\); the remaining residue class gives the third
identity.  This comparison is legitimate in the Laurent-polynomial ring
\(\IntegerNumbers[q,q^{-1}]\), where the three residue classes modulo \(3\)
form a direct sum.

Putting \(x=1\) in the last two identities gives \(B_n(1)=C_n(1)\).  Since
\(P_n(1)=0\), equation \eqref{eq:qg-borwein-dissection} at \(q=1\) gives
\(A_n(1)=B_n(1)+C_n(1)=2B_n(1)\).
\end{proof}

\begin{proposition}[Exact stabilization at both reciprocal ends]
\label{qg:prop-borwein-tail-stabilization}
Write
\[
 a_{n,r}=[x^r]A_n(x),
 \qquad
 b_{n,r}=[x^r]B_n(x),
 \qquad
 c_{n,r}=[x^r]C_n(x).
\]
For every fixed \(d\ge0\) and every \(n\ge d+1\), the lower-end
coefficients have already stabilized:
\[
 a_{n,d}=a_{d+1,d},
 \qquad
 b_{n,d}=b_{d+1,d},
 \qquad
 c_{n,d}=c_{d+1,d}.
\]
The coefficients at the same fixed distance from the reciprocal ends satisfy
the exact formulas
\begin{align}
 a_{n,n^2-d}&=a_{d+1,d},\\
 b_{n,n^2-1-d}&=c_{d+1,d},\\
 c_{n,n^2-1-d}&=b_{d+1,d}.
\label{eq:qg-borwein-tail-stabilization}
\end{align}
Thus every fixed-distance coefficient at either end of each polynomial is
eventually constant.  No coefficientwise nonnegativity assertion is needed.
\end{proposition}

\begin{proof}
Fix \(N\le n\).  Factoring the product at \(N\) gives
\[
 P_n(q)=P_N(q)
 \prod_{j=N+1}^{n}(1-q^{3j-2})(1-q^{3j-1}).
\]
For \(j>N\), both nonconstant exponents in the \(j\)th new block are at least
\(3N+1\).  Hence the entire tail product is congruent to \(1\) modulo
\(q^{3N+1}\), and
\begin{equation}
 P_n(q)\equiv P_N(q)\pmod{q^{3N+1}}.
\label{eq:qg-borwein-finite-product-stabilization}
\end{equation}
Take \(N=d+1\).  The exponents \(3d\), \(3d+1\), and \(3d+2\) are all less
than \(3N+1=3d+4\), so their coefficients in \(P_n\) and \(P_{d+1}\) agree.
By \eqref{eq:qg-borwein-dissection}, those three coefficients are respectively
\(a_{n,d}\), \(-b_{n,d}\), and \(-c_{n,d}\).  This proves the three lower-end
stabilization identities.

The reciprocity identities in \Cref{qg:prop-borwein-reciprocity} give, for
every \(n\),
\[
 a_{n,n^2-d}=a_{n,d},
 \qquad
 b_{n,n^2-1-d}=c_{n,d},
 \qquad
 c_{n,n^2-1-d}=b_{n,d}.
\]
Substituting the already proved lower-end values yields
\eqref{eq:qg-borwein-tail-stabilization} and completes the proof.
\end{proof}

\paragraph{Historical status of the sign question.}
The donor manuscript stated coefficientwise nonnegativity of
\(A_n,B_n,C_n\) as a Borwein conjecture.  That assertion is not retained here
as a conjecture: the relevant Borwein positivity conjectures are historically
resolved in the literature.  Their external proofs are not reproduced and
are not used in this chapter.  The unconditional conclusion needed here is
the sign-free stabilization theorem above, proved directly from finite-product
\(q\)-adic stabilization and reciprocity.

\section{Precisely delimited open programs}

The following items are requests for constructions or classifications, not
assertions presented as true.  They are therefore recorded as prose open
problems rather than as theorem-like conjectures.

\paragraph{Open problem (finite Andrews--Gordon certificates).}
\label{qg:prob-finite-andrews-gordon-certificates}
For every \(k\ge2\) and \(1\le i\le k\), construct a finite polynomial
identity whose coefficientwise limit is the corresponding Andrews--Gordon
identity, whose sum-side terms are manifestly coefficientwise nonnegative,
and whose verification reduces to Gaussian recurrences together with a
bounded collection of explicit telescoping-certificate types independent of
the truncation order.  Seek a canonical truncation for which the product side
has a direct lattice-path or cylindric-partition interpretation.

\paragraph{Open problem (even-modulus lowering).}
\label{qg:prob-even-modulus-lowering}
Construct a parity-sensitive parameter-lowering identity for Bailey pairs
that produces the even-modulus Bressoud family from finite algebra, with
modulus \(14\) as the first nontrivial test.  The parity restriction should be
visible in the finite kernel itself, and the proof should separate the
terminating identity, coefficientwise limit, and analytic specialization.

\paragraph{Open problem (total positivity of the Bailey kernel).}
\label{qg:prob-bailey-kernel-total-positivity}
For \(0<a,q<1\), consider the lower-triangular kernel
\[
 M_{n,r}(a,q)
 =\frac{1}{(q;q)_{n-r}(aq;q)_{n+r}}
 \qquad(0\le r\le n).
\]
Determine whether every minor of \(M\) is nonnegative (equivalently, whether
this holds after any positive diagonal normalization).  If total positivity
fails, determine the first obstructing minor and the largest order of total
positivity that holds uniformly in \(a,q\).  Either outcome should explain
which variation-diminishing or coefficient-sign consequences survive for
Bailey transforms.

\paragraph{Open problem (root-of-unity asymptotics).}
\label{qg:prob-root-of-unity-asymptotics}
Let \(\zeta\) be a primitive root of unity and approach it through
\(q=\zeta\EulerE^{-t}\), \(t\to0^+\), with all logarithmic branches stated
explicitly.  Develop uniform asymptotic expansions for q-Pochhammer products,
Bailey sums, and their locally defined inverse branches on compact parameter
sets avoiding moving poles and zeros.  Separate the dilogarithmic exponential,
algebraic prefactor, and exponentially small sectors, and determine how the
resulting Stokes data governs the singularities of the inverse q-analogs.
